\section{Repeated masks and smooth annuli}\label{sec:smooth}

The fixed-array dissipation from Section~\ref{sec:heat} controls smooth
annuli as long as we also pay for changes in the dual array. We now
allow up to $n$ disjoint intervals at every output point. The useful
feature of the jump estimate is its dependence on the dual norm:
every mixed jump term contains that norm twice. A rescaling will then
turn the linear switch count into a square-root loss.

Define the odd windows and their scale integrals by
\begin{equation}\label{eq:smooth-kernels}
 k_s(t)=\frac{t}{\sqrt s}g_s(t),\qquad
 I_{\varepsilon,R}(t)=\int_{\varepsilon^2}^{R^2}
       (k_s*k_s*k_s)(t)\,\frac{\dd s}{s}.
\end{equation}

\begin{proposition}[Smooth annuli with repeated choices]
\label{prop:smooth-switches}
Let $G_v\in C_c^\infty(\R^2)$ be complex. At each output pair $(a,b)$,
make a step choice of at most $n\ge1$ disjoint radius intervals
$(\varepsilon_j,R_j)$ and coefficients $|\beta_j|\le1$. For
\[
 \mathcal K_{a,b}(t)=\sum_j\beta_j I_{\varepsilon_j,R_j}(t)
\]
one has
\[
 |\Lambda_{\mathcal K}(G_0,G_1,G_2)|
 \le C\sqrt n\prod_{v=0}^2\|G_v\|_3.
\]
The constant is independent of the step partition, the menu of choices,
and all endpoints.
\end{proposition}

\subsection{The cost of the switches}

Work first on the finite grid of Section~\ref{sec:heat}. Make the stated
interval and coefficient choices separately for each pair $(i,j)$ and put
\begin{equation}\label{eq:repeated-mask}
 A_0^s(i,j)=A_0(i,j)\sum_{\ell}\beta_\ell
                  \mathbf1_{\{\varepsilon_\ell^2<s<R_\ell^2\}}.
\end{equation}
The sum runs through the list chosen at this particular index pair.
Keep $A_1,A_2$ fixed. Choose a compact positive scale interval
$[s_-,s_+]$ containing all switch scales in its interior, and take
$h^2\le s_-$. By disjointness,
\begin{equation}\label{eq:mask-switch-count}
 |A_0^s(i,j)|\le |A_0(i,j)|,
 \qquad
 \sum_\tau|A_0^{\tau+}(i,j)-A_0^{\tau-}(i,j)|
 \le2n|A_0(i,j)|.
\end{equation}
The sum is over the distinct switch scales, and one-sided values are
used. Both inequalities hold for complex coefficients, including
simultaneous openings and closings.

\begin{lemma}[Repeated-mask dissipation]\label{lem:repeated-dissipation}
Use the masked arrays above and variances $\sigma_v=\alpha_vs$, where
$\boldsymbol\alpha$ is any permutation of $(1,11/10,11/10)$.
With $n_v$ denoting the norms of the original arrays,
\begin{equation}\label{eq:repeated-dissipation}
 \int_{s_-}^{s_+}\int_{\Pi_0}\sum_vV_v\dd\pi_0\dd s
 \le C\left(\sum_vn_v^3+n\bigl(n_0^3+n_0^2(n_1+n_2)\bigr)\right).
\end{equation}
\end{lemma}
\begin{proof}
Between switch scales the arrays are fixed, so
\eqref{eq:heat-strict} applies. It remains to bound the energy jumps.
At a switch $\tau$, freeze the weights and set
$\delta A(i,j)=A_0^{\tau+}(i,j)-A_0^{\tau-}(i,j)$.
Interpolate all changed entries linearly. Their moduli stay bounded
by $|A_0|$ by convexity. For fixed centers and $(i,j)$ let
\begin{align*}
 X^2&=\sum_m|A_0(i,m)|^2w_1(m),&
 Y^2&=\sum_m|A_0(m,j)|^2w_0(m),\\
 H^2&=\sum_m|A_2(m,i)|^2w_2(m),&
 Q^2&=\sum_m|A_1(j,m)|^2w_2(m).
\end{align*}
These quantities use the original arrays, and square roots are
nonnegative. In $R_0=[W_0\ W_2^*]$, the changing entry has row norm
at most $\sqrt{w_0(i)}(X+H)$, column norm at most
$\sqrt{w_1(j)}Y$, and derivative modulus
$\sqrt{w_0(i)w_1(j)}|\delta A(i,j)|$.
The adjoint entry in $R_1=[W_1\ W_0^*]$ instead has row and column
bounds $\sqrt{w_1(j)}(Y+Q)$ and $\sqrt{w_0(i)}X$.
The star $R_2$ is unchanged. Lemma~\ref{lem:rank-change} therefore gives
\begin{equation}\label{eq:switch-entry-bound}
 |\Delta e_0|+|\Delta e_1|
 \le C\sum_{i,j}w_0(i)w_1(j)|\delta A(i,j)|
                    (2XY+HY+QX).
\end{equation}

The measure $h^{-2}w_0(i)w_1(j)\dd\pi_0$ is a probability measure.
Under it, $p_0,p_1$ are independent normals with means $u_i,u_j$
and variances $\sigma_0,\sigma_1$, while $p_2=-p_0-p_1$.
Adding Gaussian variances gives
\begin{align*}
 \mathbb E w_0(m)&=h g_{2\sigma_0}(u_m-u_i),\\
 \mathbb E w_1(m)&=h g_{2\sigma_1}(u_m-u_j),\\
 \mathbb E w_2(m)&=h g_{\sigma_0+\sigma_1+\sigma_2}(u_m+u_i+u_j).
\end{align*}
Extend all arrays by zero to $\mathbb Z^2$. Lemma~\ref{lem:line-maximal}
bounds $\mathbb E X^2,\mathbb E Y^2,\mathbb E H^2,\mathbb E Q^2$
by constant multiples of, respectively,
\begin{align*}
 X_*^2(i,j)&=M_{\mathbb Z}(|A_0(i,\cdot)|^2)(j),&
 Y_*^2(i,j)&=M_{\mathbb Z}(|A_0(\cdot,j)|^2)(i),\\
 H_*^2(i,j)&=M_{\mathbb Z}(|A_2(\cdot,i)|^2)(-i-j),&
 Q_*^2(i,j)&=M_{\mathbb Z}(|A_1(j,\cdot)|^2)(-i-j).
\end{align*}
All these majorants are independent of the switch scale.
Cauchy--Schwarz under the probability measure controls the three
products in \eqref{eq:switch-entry-bound}. Their square-root majorants
have $L^3(\mathbb Z^2,h^2)$ norms at most
$Cn_0,Cn_0,Cn_2,Cn_1$. This follows from the $L^{3/2}$ line maximal
inequality, applied on slices; the changes $j\mapsto-i-j$ or
$i\mapsto-i-j$ are bijections on the relevant slices.

Integrate \eqref{eq:switch-entry-bound}, sum the jumps using
\eqref{eq:mask-switch-count}, and apply H\"older on the grid. The result is
\[
 \sum_\tau\left|\Delta\sum_vJ_v(0,\tau)\right|
 \le Cn\bigl(n_0^3+n_0^2(n_1+n_2)\bigr).
\]
Finally integrate \eqref{eq:heat-strict} between consecutive switches.
The resulting sum is bounded by the initial energy plus the sum of
absolute jumps, after discarding the nonnegative terminal energy.
Lemma~\ref{lem:initial-energy} bounds the initial energy by
$C\sum_vn_v^3$. This proves \eqref{eq:repeated-dissipation}.
\end{proof}

\subsection{From the energy to the smooth kernel}

\begin{proof}[Proof of Proposition~\ref{prop:smooth-switches}]
Put $\lambda=11/10$ and use base variances $\lambda s$ at all vertices.
In Lemma~\ref{lem:mixed-trace}, insert the diagonal matrix with entries
$k_s(u_i-p_v)/g_{\lambda s}(u_i-p_v)$ at vertex $v$.
It is Hermitian and uniformly bounded: the polynomial factor in $k_s$
is absorbed by the extra Gaussian decay relative to $g_{\lambda s}$.

For the cost at vertex $v$, narrow just its row variance to $s$.
Write $T_v^{\rm b}$ for the base Gram and $T_v^{\rm n}$ for the
narrowed-row Gram. Directly from the densities,
\[
 (R_v^{\rm b})^*D_vR_v^{\rm b}=\sqrt s\,U_v^{\rm n},
 \qquad T_v^{\rm n}\le\sqrt\lambda\,T_v^{\rm b}.
\]
Invertible row scaling preserves the support. Lemma~\ref{lem:order}
therefore bounds the base cost by $C sV_v^{\rm n}$.
The narrowed variances are one of the permutations allowed in
Lemma~\ref{lem:repeated-dissipation}.

Integrate the mixed trace inequality over base centers in $\Pi_0$ and
over $s$ with measure $\dd s/s$. Its right side is bounded by
\eqref{eq:repeated-dissipation}. On the left, the density at each
vertex is now $h k_s(u_i-p_v)$. Integration over the center plane
convolves the three windows, giving
\[
 h^3\sum_{i,j,l} A_0(i,j)A_1(j,l)A_2(l,i)
       \sum_{\text{chosen at }(i,j)}\beta\,
       I_{\varepsilon,R}(u_i+u_j+u_l).
\]
All signed integrals here are absolutely convergent for fixed arrays
on the compact positive scale interval.

If the original norms are nonzero, apply this bound to arrays normalized
to have norms $n^{-1/2},1,1$. Its right side is bounded by an absolute
constant, since $n(n^{-3/2}+2n^{-1})$ is bounded. Rescaling gives
$C\sqrt n\prod_vn_v$. The zero-norm case is immediate.
The Riemann passage described in Section~\ref{sec:preliminaries}
now proves the proposition. The finitely many endpoints are fixed
before the mesh tends to zero, so its scale restriction causes no loss.
\end{proof}

\subsection{The hard endpoint errors}

The smooth annuli differ from the hard Hilbert annuli by a difference
of two localized odd kernels. We record the exact identity because
the next two sections must estimate their sum, rather than pay a
constant separately for each of the up to $2n$ endpoints.

Put
\[
 \psi=k_1*k_1*k_1=-g_3''',\qquad
 P(v)=2\int_0^v\psi(z)\dd z\quad(v\ge0),\qquad
 c_0=P(\infty)=2g_3''(0)\ne0.
\]
The function $P$ has a smooth even extension and vanishes to second
order at zero. For $\rho>0$ define
\begin{equation}\label{eq:endpoint-kernel}
 E_\rho(t)=\frac{P(|t|/\rho)-c_0\mathbf1_{\{|t|>\rho\}}}{t},
\end{equation}
with the continuous value zero at $t=0$.
Gaussian scaling and the substitution $z=|t|/\sqrt s$ in
\eqref{eq:smooth-kernels} yield, away from endpoint values,
\begin{equation}\label{eq:hard-smooth}
 I_{\varepsilon,R}(t)
 =\frac{P(|t|/\varepsilon)-P(|t|/R)}{t}
 =\frac{c_0}{t}\mathbf1_{\{\varepsilon<|t|<R\}}
       +E_\varepsilon(t)-E_R(t).
\end{equation}
Here $k_1=-g_1'$, so convolution gives $\psi=-g_3'''$; integrating
this derivative gives the stated value of $c_0$.
Each $E_\rho$ is an odd dilate of a kernel with Gaussian decay and
two jumps. The next section proves a frequency-block estimate that
can accommodate many such jumps without a linear loss in their number.
